\ifwithsols
\else
\clearpage
\fi

\subsection*{سؤال 3 امتحان 899371 سنة 2024}
\addcontentsline{toc}{subsection}{سؤال 3 امتحان 899371 سنة 2024}

\noindent
\makebox[\textwidth][c]{\includegraphics[width=\paperwidth,keepaspectratio]{ ../../../bagrut_questions/basics/arrays_2024_899371_3.jpg }%
}%

\begin{boxHint}
    \begin{itemize}
        \item ستحتاج إلى استخدام \texttt{Random} لتوليد أعداد عشوائية.
        \item استخدم حلقة \texttt{while}.
        \item تقديم العدّاد يحصل فقط عند نجاح عملية التبديل (أي عندما لا يكونا العنصرين المختارين نفسهما).
    \end{itemize}
\end{boxHint}

\ifwithsols
\begin{boxSolution}[سؤال 3 امتحان 899371 سنة 2024]

\begin{boxCode}
\begin{english}
\begin{minted}{csharp}
public static void Shuffle(int[] arr)
{
    Random rnd = new Random();
    int cnt = 0;
    while(cnt < 30)
    {
        int index1 = rnd.Next(arr.Length);
        int index2 = rnd.Next(arr.Length);
        if(index1 != index2)
        {
            int temp = arr[index1];
            arr[index1] = arr[index2];
            arr[index2] = temp;
            cnt++;
        }
    }
}
\end{minted}
\end{english}
\end{boxCode}

\end{boxSolution}
\clearpage
\fi
